\section{Related Work and Scope Boundaries}\label{sec:related}
% ============================================================
%  Section II Budget: ~0.45 page (~380--420 words, 3 compact paragraphs)
%  Note: If targeting a 4-page IEEE Letter, fold these 3 paragraphs into
%  Section I-P2 without losing citations or verification-tier boundaries.
% ============================================================

\outlinepara{II-P1: Cross-Technology Emulation \& Heterogeneous Broadcast Control}{140--155 words}{%
  \begin{itemize}
    \item \textbf{PHY-level emulation:} BlueBee~\cite{jiang2017bluebee} (extended in~\cite{jiang2019boosting}) maps \stdfour{} chips onto \SI{1}{Mbps} \ble{} GFSK payload bits ($4\times$ payload expansion), demonstrating dual-standard frame compliance and light-bulb control on commodity radios. Whereas BlueBee focuses on PHY bit-mapping and data/legacy-channel frame delivery, we examine co-locating independent \ble{} and \stdfour{} application records plus shared slot-timing metadata inside a single \ble{}~5 \texttt{AUX\_ADV\_IND} PDU under explicit retry and waveform-budget constraints.
    \item \textbf{Packet/OFDMA multi-protocol broadcast:} Brunner et al.~\cite{brunner2021leveraging} build gateway-free smart-home broadcast control over packet-level X-Burst scheduling alongside native stacks, and WiCast~\cite{xia2025wicast} synthesizes composite IEEE~802.11ax OFDMA resource units to address heterogeneous receivers in parallel. Our scope differs in transmitter constraint: a single \SI{2}{MHz} \ble{}~5 advertising radio without Wi-Fi OFDMA or packet-length energy modulation.
  \end{itemize}
  \claimboundary{Cite \cite{jiang2019boosting,xia2025wicast} strictly from verified bibliographic/abstract scope; do not assert unread internal section omissions.}
}

\outlinepara{II-P2: Cross-Technology Timing vs.\ Co-Packaged Relative Reference}{130--145 words}{%
  \begin{itemize}
    \item Dedicated CTC synchronization protocols achieve sub-millisecond or microsecond clock alignment across heterogeneous devices---for example, X-Sync~\cite{hofmann2021xsync} uses bidirectional energy bursts and RSS binary search between COTS \ble{} and \stdfour{} nodes, while Crocs~\cite{yu2018crocs} and TimeBee~\cite{gao2023timebee} address Wi-Fi-to-ZigBee time distribution.
    \item Clarify our distinct objective: we do \emph{not} propose a new cross-device clock synchronization algorithm or rank synchronization accuracy against X-Sync~\cite{hofmann2021xsync}. Instead, \proto{} embeds unidirectional slot-reference metadata within the co-packaged control stream and evaluates local IQ sample-domain held-out prediction and unilateral loss/reacquisition state transitions on a fixed secondary channel.
  \end{itemize}
}

\outlinepara{II-P3: BLE~5 Extended Advertising \& Explicit Non-Goals}{100--120 words}{%
  \begin{itemize}
    \item \ble{}~5.1 extended advertising~\cite{bluetooth_core51} decouples primary pointers (\texttt{ADV\_EXT\_IND} on ch37--39) from secondary payloads (\texttt{AUX\_ADV\_IND}, up to \SI{255}{B}), removing the \SI{31}{B} legacy ceiling that forces multi-frame fragmentation (\fragproto{}).
    \item State explicit scope and attribution boundaries upfront: confirmatory runs use a fixed secondary channel (avoiding conflation with \texttt{AuxPtr} channel-hopping transients), where \modejc{} and \modesc{} share \ble{} ch3 (\SI{2410}{MHz}) and \stdfour{} ch12 (\SI{2410}{MHz}) for protocol-identifiable comparison, whereas \modelc{} transmits \ble{} on legacy ch39 (\SI{2480}{MHz}) so \modejc{}--\modelc{} resolution differences apply strictly to the tested encapsulation+channel combination. Waveform duration (\si{\micro\second}) is a submitted design proxy rather than power-meter RF energy, and independent SDR receivers do not claim commercial stack compatibility or physical actuator synchronization.
  \end{itemize}
}
